\subsection{The ae version of the main theorem}\label{sec:main-theorem-ae}

An adjunction of the same shape as the main theorem holds in the category
obtained by identifying morphisms up to null functions
(Definition~\ref{def:int-ae}).

\begin{lemma}\label{lem:daniell-extension-preserves-ae}
  Let $(X,L,I)$ and $(Y,M,J)$ be integration systems and let
  \[
    F,G:(X,L,I)\to(Y,M,J)
  \]
  be morphisms of integration systems. Then
  \[
    F=_{\mathrm{ae}}G
    \quad\Longrightarrow\quad
    F^1=_{\mathrm{ae}}G^1 ,
  \]
  where $F^1,G^1$ are the morphisms
  \[
    (X,L^1,I^1)\to(Y,M^1,J^1)
  \]
  induced by the Daniell extension.
\end{lemma}
\begin{proof}
  Take $f\in L^1$ and take the sequence of functions $\{s_N\}\subset L$ of
  Lemma~\ref{lem:l1-approx-by-source}.
  Since $F^1,G^1$ are morphisms of Riesz systems and of integration systems and
  $J^1$ extends $J$,
  \[
    \begin{aligned}
      &J^1\bigl(|(F^1)^*f-(G^1)^*f|\bigr)\\
      &\quad\le
      J^1\bigl(|(F^1)^*(f-s_N)|\bigr)
      +J^1\bigl(|F^*s_N-G^*s_N|\bigr)
      +J^1\bigl(|(G^1)^*(s_N-f)|\bigr)\\
      &\quad=
      I^1(|f-s_N|)
      +J\bigl(|F^*s_N-G^*s_N|\bigr)
      +I^1(|s_N-f|)\\
      &\quad=2I^1(|f-s_N|).
    \end{aligned}
  \]
  The last equality used $F=_{\mathrm{ae}}G$ and $s_N\in L$.
  Letting $N\to\infty$, the left-hand side is $0$. As $f\in L^1$ was arbitrary,
  $F^1=_{\mathrm{ae}}G^1$.
\end{proof}

\begin{corollary}\label{cor:daniell-extension-ae-functor}
  The assignment sending an object to its Daniell extension and a morphism to
  \[
    [F]\longmapsto[F^1]
  \]
  defines a functor
  \[
    \mathcal{D}_{\mathrm{ae}}:
    \mathbf{Int}_{\mathrm{ae}}
    \to
    \mathbf{DInt}_{\mathrm{ae}} .
  \]
\end{corollary}
\begin{proof}
  By Lemma~\ref{lem:daniell-extension-preserves-ae} the assignment on morphisms
  does not depend on the choice of representative.
  Moreover, that the Daniell extension is a Daniell closed integration system
  and that identities and composition of morphisms are preserved were shown up
  to Subsection~\ref{sec:main-theorem}.
\end{proof}

\begin{definition}
  We denote the inclusion functor of the full subcategory by
  \[
    \iota_{D,\mathrm{ae}}:
    \mathbf{DInt}_{\mathrm{ae}}
    \to
    \mathbf{Int}_{\mathrm{ae}} .
  \]
\end{definition}

\begin{theorem}[ae version of the main theorem]\label{thm:main-ae}
  The functor
  \[
    \mathcal{D}_{\mathrm{ae}}:
    \mathbf{Int}_{\mathrm{ae}}
    \to
    \mathbf{DInt}_{\mathrm{ae}}
  \]
  defined by the Daniell extension is left adjoint to the inclusion functor
  $\iota_{D,\mathrm{ae}}$.
\end{theorem}
\begin{proof}
  Let $A=(X,L,I)$ be an integration system and write its Daniell extension as
  $\mathcal{D}_{\mathrm{ae}}A=(X,L^1,I^1)$.
  We show that the ae-equivalence class
  \[
    [i_A]:A\to\iota_{D,\mathrm{ae}}\mathcal{D}_{\mathrm{ae}}A
  \]
  of the morphism of integration systems
  \[
    i_A:A\to\iota_{D}\mathcal{D}A
  \]
  constructed in Subsection~\ref{sec:main-theorem} is a universal morphism from
  $A$ to $\iota_{D,\mathrm{ae}}$.

  Let $B=(Y,M,J)$ be a Daniell closed integration system and let
  \[
    [F]:A\to\iota_{D,\mathrm{ae}}B
  \]
  be a morphism. Choose a representative $F$. By the proof of
  Theorem~\ref{thm:main-theorem} and $M^1=M$, there is a morphism
  \[
    F^1:\mathcal{D}_{\mathrm{ae}}A\to B
  \]
  induced by the Daniell extension such that, at the level of representatives,
  \[
    F=\iota_{D}(F^1)\circ i_A .
  \]
  Hence
  \[
    [F]=\iota_{D,\mathrm{ae}}([F^1])\circ[i_A].
  \]

  Next, suppose a morphism $[H]:\mathcal{D}_{\mathrm{ae}}A\to B$ satisfies
  \[
    [F]=\iota_{D,\mathrm{ae}}([H])\circ[i_A] .
  \]
  This means that, for representatives,
  \[
    F=_{\mathrm{ae}}H\circ i_A .
  \]
  Applying Lemma~\ref{lem:daniell-extension-preserves-ae} to these two
  morphisms,
  \[
    F^1=_{\mathrm{ae}}(H\circ i_A)^1 .
  \]
  The Daniell extension does not change the map of base spaces underlying a
  morphism, and $B$ is Daniell closed, so $(H\circ i_A)^1=H$. Hence
  \[
    [F^1]=[H].
  \]
  Therefore the ae-equivalence class of the factorizing morphism is unique and
  $[i_A]$ is a universal morphism.

  Since $A$ was arbitrary, $\iota_{D,\mathrm{ae}}$ has a left adjoint.
  Moreover, for a morphism of integration systems $T:A\to A'$ the class $[T^1]$
  satisfies
  \[
    \iota_{D,\mathrm{ae}}([T^1])\circ[i_A]
    =[i_{A'}]\circ[T] ,
  \]
  so the action on morphisms of the left adjoint obtained from the universal
  property agrees with $\mathcal{D}_{\mathrm{ae}}$.
\end{proof}

\begin{corollary}
  The functor
  \[
    \mathcal{D}_{M,\mathrm{ae}}:
    \mathbf{MInt}_{\mathrm{ae}}
    \to
    \mathbf{DInt}_{\mathrm{ae}}
  \]
  obtained by restricting $\mathcal{D}_{\mathrm{ae}}$ to
  $\mathbf{MInt}_{\mathrm{ae}}$ is left adjoint to the inclusion functor
  \[
    \iota_{M,\mathrm{ae}}:
    \mathbf{DInt}_{\mathrm{ae}}
    \to
    \mathbf{MInt}_{\mathrm{ae}} .
  \]
\end{corollary}
\begin{proof}
  $\mathbf{MInt}_{\mathrm{ae}}$ is a full subcategory of
  $\mathbf{Int}_{\mathrm{ae}}$, and $\mathbf{DInt}_{\mathrm{ae}}$ is a full
  subcategory of it.
  It therefore suffices to restrict the adjunction isomorphism of
  Theorem~\ref{thm:main-ae} to the objects of
  $\mathbf{MInt}_{\mathrm{ae}}$.
\end{proof}

\begin{remark}\label{rem:ae-mct-counterexample-fails}
  Dividing morphisms by ae-equivalence does not weaken the membership condition
  for functions in the target of the factorization.
  Accordingly, in the theorem above as in the main theorem, Daniell closedness is
  assumed of the target of the factorization.

  However, \textbf{the finite counterexample of
  Remark~\ref{rem:D-I-not-left-adjoint} cannot be used as it stands as grounds
  for the claim that Daniell closedness cannot be dropped.}
  In that example only $\tau(0)=0$ has any effect, so all $5$ morphisms and all
  $3$ morphisms are ae-equivalent to one another, and in the ae-category both
  $\operatorname{Hom}$ sets collapse to a single point.
  A counterexample for the ae version is given in the next subsubsection
  (Corollary~\ref{cor:ae-completeness-needed}).

  This point contrasts with Section~\ref{chap:operator}.
  There, once morphisms are replaced by operators, the extension is determined
  only up to null functions, and dividing the codomain by null functions becomes
  the natural thing to do.
  Even there, however, we do not prove that ``it is undetermined unless one
  divides'' (Problem~\ref{prob:no-lift}).
\end{remark}

\subsubsection{Daniell closedness cannot be dropped in the ae version either}

As stated in Remark~\ref{rem:ae-mct-counterexample-fails}, the finite
counterexample of Remark~\ref{rem:D-I-not-left-adjoint} has no effect in the
ae-category.
A counterexample for the ae version is given by the pair consisting of the
integration system of Borel measurable functions and the integration system of
Lebesgue integrable functions.

\begin{proposition}\label{prop:no-morphism-lebesgue-to-borel}
  Let $X=[0,1)$ and put
  \begin{align*}
    P&:=\bigl(X,\ \{f:X\to\mathbb{R}\text{ Borel measurable and Lebesgue integrable}\},\ \text{the Lebesgue integral}\bigr),\\
    Q&:=\bigl(X,\ \{f:X\to\mathbb{R}\text{ Lebesgue integrable}\},\ \text{the Lebesgue integral}\bigr).
  \end{align*}
  Then
  \[
    \operatorname{Hom}_{\mathbf{Int}}(Q,P)=\emptyset.
  \]
\end{proposition}
\begin{proof}
  Suppose a morphism $F:Q\to P$ exists and let $\tau:X\to X$ be its map of base
  spaces.
  By definition, for every Lebesgue integrable $h:X\to\mathbb{R}$ the function
  $h\circ\tau$ is Borel measurable and Lebesgue integrable, and
  \[
    \int_X h\circ\tau\,d\lambda=\int_X h\,d\lambda
  \]
  holds, where $\lambda$ is Lebesgue measure.

  \textbf{$\tau$ is Borel measurable and measure-preserving.}
  For a Borel set $B\subset X$ the function $1_B$ is Lebesgue integrable, and
  since $1_B\circ\tau=1_{\tau^{-1}(B)}$ is Borel measurable, the set
  $\tau^{-1}(B)$ is Borel.
  Moreover preservation of the integral gives
  $\lambda(\tau^{-1}(B))=\lambda(B)$.

  \textbf{The image of $\tau$ is uncountable.}
  If $\tau(X)$ were contained in a Borel null set $N_0$, then
  $\tau^{-1}(X\setminus N_0)=\emptyset$ while
  $\lambda(X\setminus N_0)=1$, contradicting preservation of measure.
  A countable set is contained in a Borel null set, so $\tau(X)$ is
  uncountable.

  \textbf{The image of $\tau$ contains a Lebesgue null set of cardinality
  $\mathfrak{c}$.}
  Since $\tau$ is Borel measurable, its image $\tau(X)$ is an analytic set.
  An uncountable analytic set contains a nonempty perfect set (the perfect set
  property of analytic sets; see \cite{Kechris}).
  Carrying out a Cantor-type construction inside such a perfect set $K$, with
  all $2^n$ closed intervals at stage $n$ of diameter at most $4^{-n}$, the
  resulting intersection $P_0\subset K$ is a set of cardinality $\mathfrak{c}$,
  and $\lambda(P_0)\le2^n\cdot4^{-n}=2^{-n}$ for every $n$, so
  $\lambda(P_0)=0$.

  \textbf{Contradiction.}
  Take an arbitrary $S\subset P_0$. As a subset of a Lebesgue null set, $S$ is
  Lebesgue measurable, and $1_S$ is Lebesgue integrable.
  Hence $1_S\circ\tau=1_{\tau^{-1}(S)}$ is Borel measurable, that is,
  $\tau^{-1}(S)$ is a Borel set.

  The assignment $S\mapsto\tau^{-1}(S)$ is injective on the subsets of $P_0$.
  Indeed, if $S\ne S'$ then, without loss of generality, we may take
  $p\in S\setminus S'$.
  Since $p\in P_0\subset\tau(X)$ there is an $x\in X$ with $\tau(x)=p$, and
  $x\in\tau^{-1}(S)\setminus\tau^{-1}(S')$.

  We thus obtain an injection from $\mathcal{P}(P_0)$ into the set of all Borel
  subsets of $X$.
  But $|\mathcal{P}(P_0)|=2^{\mathfrak{c}}$ while $X$ has only $\mathfrak{c}$
  Borel sets (the same fact as used in the proof of
  Corollary~\ref{cor:lebesgue-reflections-differ}).
  This is a contradiction.
\end{proof}

\begin{corollary}\label{cor:ae-completeness-needed}
  The functor
  \[
    \iota_{M,\mathrm{ae}}\circ\mathcal{D}_{\mathrm{ae}}:
    \mathbf{Int}_{\mathrm{ae}}\longrightarrow\mathbf{MInt}_{\mathrm{ae}}
  \]
  defined by the Daniell extension is not left adjoint to the inclusion functor
  $\iota_{\mathrm{ae}}:\mathbf{MInt}_{\mathrm{ae}}\to\mathbf{Int}_{\mathrm{ae}}$
  of the full subcategory.
  That is, \textbf{imposing Daniell closedness on the codomain of the
  factorization cannot be dispensed with in the ae version either}.
\end{corollary}
\begin{proof}
  Taking $A$ to be the integration system of
  Example~\ref{ex:step-functions}, by
  Corollary~\ref{cor:lebesgue-reflections-differ} the object $P$ of
  Proposition~\ref{prop:no-morphism-lebesgue-to-borel} is $A^{\mathrm{M}}$ and
  $Q$ is $\mathcal{D}A$.
  Here $P$ is an object of $\mathbf{MInt}$
  (Theorem~\ref{thm:mint-inclusion-has-left-adjoint}) and is not Daniell closed
  (the same corollary).
  Moreover $\mathcal{D}P=\mathcal{D}(A^{\mathrm{M}})=\mathcal{D}A=Q$ by
  Lemma~\ref{lem:two-reflections-same-daniell}.

  The morphisms of $\mathbf{Int}_{\mathrm{ae}}$ are the ae-equivalence classes of
  morphisms of $\mathbf{Int}$, so
  Proposition~\ref{prop:no-morphism-lebesgue-to-borel} gives
  \[
    \operatorname{Hom}_{\mathbf{MInt}_{\mathrm{ae}}}(\mathcal{D}_{\mathrm{ae}}P,\,P)=\emptyset .
  \]
  On the other hand
  $\operatorname{Hom}_{\mathbf{Int}_{\mathrm{ae}}}(P,\iota_{\mathrm{ae}}P)$
  contains the ae-equivalence class of the identity morphism and is therefore
  nonempty.
  There is no bijection between the empty set and a nonempty set, so no
  adjunction isomorphism exists.
\end{proof}

\begin{remark}
  This proof is the only place in this paper that uses a result of descriptive
  set theory (the perfect set property of analytic sets).
  This contrasts with the non-ae counterexample
  (Remark~\ref{rem:D-I-not-left-adjoint}), for which a three-point set
  sufficed.
  Once morphisms are divided by ae-equivalence, finite objects can no longer
  tell them apart, and a counterexample appears only in the form of the
  difference between Borel and Lebesgue.
\end{remark}
